% P47 | 5.8 Profit and loss | Result: profit and loss under the real-world measure, and tail risk
% guide: dev_notes/paragraph_guides/P47_result_profit_and_loss_under_the_real_world_meas.txt
Under a real-world drift of $\mu=5\%$ the mean profit of the buyer is \PnLPmean, and at $\mu=10\%$ it is \PnLPmeanTen.
A drift above the risk-free rate makes the stock rise faster than under the pricing measure, so the put pays out less often than the price assumes.
The probability that the buyer loses the entire premium rises from \PnLQtot\% under the pricing measure to \PnLPtot\% at $\mu=5\%$.
The probability of exercising early falls to \PnLPpex, and the mean stopping time rises to \PnLPtau\ years.
The buyer's loss is capped at the premium, so the tail risk is on the seller's side, where the loss equals the buyer's profit.
The $99\%$ quantile of the buyer's profit is \PnLPqNinetyNine\ under the real-world measure and \PnLQqNinetyNine\ under the pricing measure, and the expected shortfall at $95\%$ is \PnLPes\ against \PnLQes, so the right tail hardly moves.
The real-world drift changes how often the put pays out, not how much it pays, which is why the mean moves while the tail does not.
A seller who hedges dynamically is insulated from the drift, as Section~\ref{sec:hedge} showed for the call, whereas an unhedged seller keeps the drift risk.
